
\begin{figure}[!t]
    \centering
    \begin{tikzpicture}[
        vertex/.style={
            circle,
            draw,
            minimum size=6mm,
            inner sep=0pt,
            line width=0.6pt
        },
        leftvertex/.style={
            vertex,
            fill=blue!18,
            draw=blue!60!black
        },
        rightvertex/.style={
            vertex,
            fill=orange!20,
            draw=orange!70!black
        },
        node distance=9mm and 25mm
    ]

        % Users
        \node[leftvertex] (u1) {};
        \node[leftvertex, below=of u1] (u2) {};
        \node[leftvertex, below=of u2] (u3) {};
        \node[leftvertex, below=of u3] (u4) {};

        % User labels
        \node[left=3mm of u1] {\small Alice};
        \node[left=3mm of u2] {\small Bob};
        \node[left=3mm of u3] {\small Carol};
        \node[left=3mm of u4] {\small David};

        % Products
        \node[rightvertex, right=of u1] (p1) {};
        \node[rightvertex, below=of p1] (p2) {};
        \node[rightvertex, below=of p2] (p3) {};
        \node[rightvertex, below=of p3] (p4) {};

        % Product labels
        \node[right=3mm of p1] {\small Laptop};
        \node[right=3mm of p2] {\small Headphones};
        \node[right=3mm of p3] {\small Keyboard};
        \node[right=3mm of p4] {\small Camera};

        % Edges forming a 3x3 biclique
        \draw (u1) -- (p1);
        \draw (u1) -- (p2);
        \draw (u1) -- (p3);

        \draw (u2) -- (p1);
        \draw (u2) -- (p2);
        \draw (u2) -- (p3);

        \draw (u3) -- (p1);
        \draw (u3) -- (p2);
        \draw (u3) -- (p3);

        % Additional edges outside the biclique
        \draw (u4) -- (p2);
        \draw (u4) -- (p4);

        % Part labels
        \node[above=3mm of u1] {\small Users};
        \node[above=3mm of p1] {\small Products};

    \end{tikzpicture}
    \caption{A bipartite user-product graph. Alice, Bob, and Carol each purchased the same three products, forming a $3\times3$ biclique, while David has a different set of purchases.}
    \label{fig:bipartite-example}
\end{figure}
